\section{The $D$ Notation}

\begin{enumerate}
  \item What is the domain and codomain of the function $f(x,y) = x^2 + y^2$?
  \item What is the domain and codomain of the derivative $Df(x,y)$?
  \item In general, for a function $f : \mathbb{R}^m \to \mathbb{R}^n$, what is the type (domain and codomain) of $Df(x)$?
  \item Why is $Df(x,y)$ considered a linear map rather than just a collection of partial derivatives?
  \item Compute the partial derivatives $\frac{\partial f}{\partial x}$ and $\frac{\partial f}{\partial y}$ for $f(x,y) = x^2 + y^2$.
  \item Write the gradient $\nabla f(x,y)$.
  \item Express $Df(x,y)$ as a row vector (Jacobian matrix).
  \item Let $h = (h_1, h_2) \in \mathbb{R}^2$. Compute $Df(x,y)(h)$.
  \item How does the expression $Df(x,y)(h)$ relate to the dot product with the gradient?
  \item Write the first-order linear approximation of $f(x,y)$ at a point $(x,y)$ using $Df(x,y)$.
  \item What is the geometric meaning of the vector $\nabla f(x,y)$?
  \item What does $Df(x,y)(h)$ represent geometrically?
  \item How does the definition of $Df(x,y)(h)$ relate to the directional derivative?
  \item Compare the expressions:
    \[
      Df(x)(h) = 2x \cdot h \quad \text{and} \quad Df(x,y)(h) = 2x \cdot h_1 + 2y \cdot h_2.
    \]
    What structural similarity do you observe?
  \item In the scalar-to-scalar case $f : \mathbb{R} \to \mathbb{R}$, why can $Df(x)$ be identified with multiplication by a number?
  \item Explain why $Df(x,y)$ can be viewed as the “best linear approximation” of $f$ near $(x,y)$.
  \item If $h = (1,0)$, what does $Df(x,y)(h)$ compute? What about $h = (0,1)$?
  \item How would you interpret $Df(x,y)$ acting on an arbitrary direction vector $h$ in terms of rates of change?
\end{enumerate}

\vspace{5mm}

\begin{center}
  \begin{tabular}{ c l l }
    Scalar $\to$ Scalar & $f : \mathbb{R} \to \mathbb{R}$     & $Df(x)(h) = f'(x)\,h$ \\
    Scalar $\to$ Vector & $f : \mathbb{R} \to \mathbb{R}^n$   & $Df(x)(h) = f'(x) \cdot h$ \\
    Vector $\to$ Scalar & $f : \mathbb{R}^m \to \mathbb{R}$   & $Df(x)(h) = \nabla f(x) \cdot h$ \\
    Vector $\to$ Vector & $f : \mathbb{R}^m \to \mathbb{R}^n$ & $Df(x)(h) = J_f(x)\,h$
  \end{tabular}
\end{center}

\vspace{5mm}

A hint: Think of $D$ as a type-lifting operator that transforms a function into a new function that returns linear maps
\[ D : ( \mathbb{R}^m \to \mathbb{R}^n ) \;\to\; ( \mathbb{R}^m \to \mathcal{L}(\mathbb{R}^m, \mathbb{R}^n) ) \]
where $\mathcal{L}(\cdot)$ = space of linear maps.

\vspace{5mm}

A hint: Think of $Df(x)$ as: ``Give me a direction $h$, and I'll tell you how $f$ changes in that direction.''
\[ Df(x)(h) \approx f(x + h) - f(x) \]
